\vspace*{2cm}
\begin{center}
{\Large \bfseries Resumen}

\vspace{2.5cm}
\end{center}

En los últimos años, la seguridad en el desarrollo de software ha
adquirido una relevancia crítica debido al aumento de vulnerabilidades y
ataques contra sistemas informáticos. Como respuesta, los compiladores
modernos han incorporado mecanismos de protección que refuerzan la
integridad y robustez de los programas generados. Sin embargo, estas
medidas introducen un coste en términos de eficiencia, afectando al
tiempo de ejecución, el uso de memoria y el rendimiento general del
software.

Este Trabajo de Fin de Grado realiza un análisis comparativo del impacto
que tienen las medidas de seguridad aplicadas durante la compilación
sobre la eficiencia de los programas generados. Para ello, se han
evaluado diversos compiladores ampliamente utilizados, como \emph{GCC}, 
\emph{Clang} y \emph{Rustc}, combinando distintas opciones de seguridad con
niveles de optimización. Durante la experimentación, se han recogido métricas
clave como el tiempo de ejecución, el uso de memoria, el tamaño del archivo resultante
y la presencia de funciones fortificadas, que reemplazan fragmentos de
código potencialmente peligrosos por versiones más seguras.

Este trabajo proporciona una visión clara y práctica sobre el coste
asociado a la seguridad en la compilación de programas, ofreciendo
recomendaciones específicas para distintos escenarios, desde
aplicaciones críticas hasta entornos donde el rendimiento es la
prioridad. Los hallazgos contribuyen a fundamentar decisiones técnicas
en proyectos donde ambos factores, seguridad y eficiencia, son
determinantes.

\clearpage

\vspace*{2cm}
\begin{center}
{\Large \bfseries Abstract}

\vspace{2.5cm}
\end{center}

In recent years, security in software development has gained critical
importance due to the rise in vulnerabilities and attacks on computer
systems. In response, modern compilers have incorporated protection
mechanisms that enhance the integrity and robustness of the generated
programs. However, these measures come at a cost in terms of efficiency,
affecting runtime, memory usage, and overall software performance.

This Bachelor's Thesis presents a comparative analysis of the impact
that security measures applied during compilation have on the efficiency
of the resulting programs. To achieve this, widely used compilers such
as \emph{GCC}, \emph{Clang} and \emph{Rustc} are evaluated, combining different
security options with various optimization levels. During
experimentation, key metrics such as execution time, memory usage,
resulting file size, and the presence of fortified functions---which
replace potentially dangerous code segments with safer versions---were
collected.

This work provides a clear and practical perspective on the cost
associated with security in program compilation, offering specific
recommendations for different scenarios, ranging from critical
applications to environments where performance is the priority. The
findings help support technical decisions in projects where both
security and efficiency are crucial factors.

\clearpage